% ------------------------------------------------------------------
%  Capitolo 9 — Il tempo alleato
%  Parte III
%  Ricerca di riferimento: ricerca 01 §2.5, §4, §5.2, §7.6, §11.1;
%    02 §9.3, §10; 03 §3.5, §3.8, §6
%  Voci bibliografiche principali: cepeda2006, cepeda2008, latimier2021,
%    rawson2013, rawson2022, kornell2009, murray2025, donoghue2021,
%    ebbinghaus1885, jenkins1924, yoo2007
%  Epigrafe: ✔ verificata sul testo (vedi ricerca 03 e registro, ricerca 00)
%  Scheda del capitolo: ricerca/06_indice-ragionato.md, capitolo 9
%
%  STATO: prima stesura (23/09/2026), rivista lo stesso giorno (intervalli, Quintiliano, riquadro docenti); S-41 chiuso.
%  Dati verificati: S-05 (Ebbinghaus 68/38), S-08 (Cepeda 2006: 839
%  confronti, 317 esperimenti, 184 articoli), S-09 (Cepeda 2008: 1.354
%  partecipanti; ottimo ≈ 20% per poche settimane, ≈ 5% per un anno),
%  S-10 (Kornell 2009: 63 su 70; 42 su 58). Verificati sull'abstract il
%  23/09/2026: Latimier et al. 2021 (29 studi; g = 0,74; crescenti vs
%  uniformi g = 0,034, n.s.) e Donoghue & Hattie 2021 (pratica
%  distribuita d = 0,85, la più alta delle dieci tecniche). Rawson et
%  al. 2013 e Murray et al. 2025 citati solo in forma qualitativa.
%  Ovidio: solo il primo emistichio è autentico (ricerca 03 §3.8, §6);
%  del contesto di Ex Ponto IV, 10 si dice solo il tema (il tempo che
%  consuma ogni cosa), senza citare altri versi non verificati.
% ------------------------------------------------------------------
\chapter{Il tempo alleato}
\label{cap:tempo}

\epigrafe[e, se possibile ogni giorno, è la cosa più efficace]{et si fieri potest cotidie, potentissimum est}{Quintiliano, Institutio oratoria XI, 2, 40}

\iniziale{E}{siliato a Tomi}, sul Mar Nero, lontano da Roma, Ovidio
scrive a un amico una lettera in versi sul tempo che
consuma ogni cosa. Tra gli esempi che porta c'è un verso che tutti
abbiamo sentito almeno una volta:
\emph{gutta cavat lapidem}, \enquote{la goccia scava la
pietra}.\footnote{Ovidio, \emph{Epistulae ex Ponto} IV, 10, 5:
\emph{gutta cavat lapidem, consumitur anulus usu}, \enquote{la goccia
scava la pietra, l'anello si consuma con l'uso}. Il seguito che spesso
si aggiunge, \emph{non vi sed saepe cadendo} (\enquote{non con la
forza, ma cadendo spesso}), non è di Ovidio: è un'aggiunta
medievale.} Ovidio parla dell'usura, di ciò che il tempo porta via. Ma
la sua immagine dice anche il contrario, e lo dice meglio di qualunque
manuale: una goccia sola non fa nulla; una secchiata d'acqua rovesciata
tutta in una volta scivola via; la stessa acqua, caduta goccia dopo
goccia per mesi, lascia un segno che resta.

Prima di andare avanti, prova a ricordare. Nel primo capitolo abbiamo
incontrato Hermann Ebbinghaus e un suo esperimento su se stesso: quante
ripetizioni gli servivano per imparare una lista di sillabe se le faceva
tutte in un giorno, e quante se le distribuiva su tre giorni? Se non
ricordi i numeri, non importa: prova comunque, e poi guarda la
risposta nella nota.\footnote{Sessantotto ripetizioni fatte tutte di
seguito, contro trentotto distribuite nei tre giorni precedenti, per lo
stesso risultato: ripetere la lista senza errori il giorno dopo
\parencite{ebbinghaus1885}.} Lo sforzo che hai appena fatto è pratica
del recupero (capitolo~\ref{cap:recupero}); il fatto che tu lo abbia
fatto qualche capitolo dopo, e non subito, è il tema di questo
capitolo.

La conclusione di Ebbinghaus era che distribuire le ripetizioni nel
tempo è \enquote{decisamente più vantaggioso} che accumularle in
un'unica sessione. Quasi la metà del lavoro, per lo stesso risultato.
Da allora la psicologia ha misurato questo fenomeno centinaia di volte,
in quasi tutti i modi possibili, e ha scoperto qualcosa che Ebbinghaus
poteva solo intuire: il vantaggio della distribuzione cresce quanto più
a lungo si vuole ricordare. Per un esame di domani, concentrare lo
studio funziona abbastanza. Per ricordare tra un anno, o per la vita,
non funziona quasi per niente. Questo capitolo spiega perché il tempo,
che di solito consideriamo il nemico della memoria, è in realtà il suo
alleato più prezioso; quando conviene tornare su ciò che hai studiato;
e come si costruisce, in pratica, un calendario di ripasso.

\section{La goccia che scava la pietra}

Gli psicologi chiamano questo fenomeno \emph{effetto della
distribuzione} (in inglese \emph{spacing effect}) e lo descrivono così:
a parità di tempo totale, studiare in più sessioni separate da
intervalli produce un ricordo a lungo termine molto migliore che
studiare in un'unica sessione. Il contrario della pratica distribuita
si chiama \emph{pratica concentrata} (\emph{massed practice}); nella
sua forma estrema, quella della vigilia dell'esame, gli studenti di
lingua inglese la chiamano \emph{cramming}, da un verbo che significa
\enquote{stipare}, come si stipano i vestiti in una valigia troppo
piccola.

È uno dei risultati più solidi e più replicati della psicologia
sperimentale. Nel 2006 Nicholas Cepeda e i suoi colleghi hanno raccolto
tutti gli studi che erano riusciti a trovare sulla distribuzione nel
tempo di compiti di memoria verbale: 184 articoli, 317 esperimenti, 839
confronti tra studio distribuito e studio concentrato
\parencite{cepeda2006}. Il vantaggio della distribuzione compariva
quasi ovunque: con liste di parole, con coppie di termini, con fatti
da ricordare, con intervalli di minuti e con intervalli di settimane. E quando, quindici
anni dopo, John Donoghue e John Hattie hanno messo a confronto le dieci
tecniche di studio più diffuse raccogliendo più di duecento studi, la
pratica distribuita è risultata quella con l'effetto medio più alto di
tutte, davanti alla stessa pratica del recupero
\parencite{donoghue2021}.

Perché dovrebbe funzionare così? La risposta migliore che abbiamo è
quella che abbiamo incontrato nel capitolo~\ref{cap:dueforze}. Ogni
ricordo ha due forze: quanto è radicato (la forza di immagazzinamento)
e quanto è accessibile in questo momento (la forza di recupero). Quando
ripeti una cosa appena studiata, la sua forza di recupero è ancora
altissima: la ripetizione non costa nulla, e proprio per questo
aggiunge poco. Quando invece torni su quella cosa dopo un giorno, o
dopo una settimana, una parte è già sbiadita. Ricordarla costa uno
sforzo, e quello sforzo -- se riesce -- la radica molto più a fondo.
Detto in modo paradossale: \emph{un po' di oblio tra una ripetizione e
l'altra rende la ripetizione successiva più efficace}. È un esempio
perfetto di quelle che Robert ed Elizabeth Bjork hanno chiamato
\emph{difficoltà desiderabili} \parencite{bjork2011}: condizioni che rendono lo studio più
faticoso oggi e il ricordo più saldo domani.

Quintiliano, che non conosceva nessuna di queste teorie, aveva visto il
fenomeno con i propri occhi. Nella frase che apre questo capitolo
dice che cosa risponderebbe a chi gli chiedesse quale sia la più
grande arte della memoria:
\enquote{l'esercizio e la fatica: imparare molte cose, riflettere su
molte cose, e, se possibile ogni giorno, è la cosa più efficace:
niente come la memoria cresce con la cura o svanisce con la
trascuratezza}.\footnote{Quintiliano, \emph{Institutio oratoria} XI,
2, 40.} Poco oltre aggiunge un'osservazione ancora più sottile: ciò che
sul momento non si riusciva a ripetere, il giorno dopo \enquote{si
ricompone}, e \enquote{quello stesso tempo che di solito è causa di
oblio rafforza la memoria}.\footnote{Quintiliano, \emph{Institutio
oratoria} XI, 2, 43. Torneremo su questo passo nel
capitolo~\ref{cap:corpo}, a proposito del sonno.} Il tempo che fa
dimenticare è lo stesso che fa ricordare. Quintiliano osserva il
fenomeno e ammette di non saperlo spiegare; non è una prova, ma è
un'intuizione che la psicologia sperimentale avrebbe confermato quasi
diciotto secoli dopo.

C'è un'ultima ragione, più semplice, per cui distribuire funziona:
torni sulle cose in giorni diversi, con uno stato d'animo diverso,
magari in un luogo diverso, dopo aver studiato altro. Il ricordo si
lega a una varietà di contesti e diventa meno dipendente da quello in
cui è nato. Chi ha studiato un argomento soltanto in una notte, nella
stessa stanza, sotto la stessa lampada, possiede un ricordo che
funziona bene in quelle condizioni e molto meno fuori di esse. L'esame,
e ancora di più la vita, non avvengono in quella stanza.

\section{La cresta dell'oblio}

Se distribuire conviene, la domanda successiva è ovvia: \emph{quanto}
distribuire? Tra una sessione di studio e la ripresa, meglio lasciar
passare un giorno, una settimana, un mese?

Per rispondere, nel 2008 lo stesso gruppo di Cepeda ha condotto uno
degli esperimenti più ambiziosi della storia di questo campo
\parencite{cepeda2008}. Più di milletrecento persone hanno imparato una
serie di fatti poco noti (per esempio, quale sia il paese europeo che
mangia più cibo piccante messicano). Poi hanno ripreso lo stesso
materiale dopo un intervallo che variava, a seconda del gruppo, da
nessun giorno a più di tre mesi. Infine sono state messe alla prova a
distanze diverse dall'ultima ripresa: una settimana, cinque settimane,
dieci settimane, quasi un anno. Incrociando gli intervalli di ripresa e
gli intervalli del test, i ricercatori hanno potuto disegnare una sorta
di mappa, e su quella mappa è comparsa una linea: una \emph{cresta}
lungo la quale il ricordo era massimo.

La cresta dice una cosa molto semplice e molto utile: \emph{l'intervallo
migliore tra due sessioni è una frazione del tempo per cui vuoi
ricordare}. Se devi ricordare qualcosa per poche settimane, conviene
riprenderlo dopo circa un quinto di quel tempo; se devi ricordarlo per
un anno, la frazione scende a circa un ventesimo, cioè da due a quattro
settimane. Per una prova tra dieci giorni, la ripresa ideale cade dopo
un paio di giorni; per un esame alla fine del semestre, dopo due o tre
settimane; per un concorso tra un anno, dopo qualche settimana. Chi
riprende tutto il giorno dopo, per un esame tra sei mesi, ha
distribuito troppo poco; chi aspetta tre mesi, per una prova tra
quindici giorni, ha semplicemente sbagliato i conti.

La stessa ricerca ha mostrato un'asimmetria che vale oro nella pratica:
\emph{intervalli un po' troppo lunghi costano poco; intervalli troppo
brevi costano molto}. Superata la cresta, il ricordo cala lentamente;
prima della cresta, precipita. Se hai un dubbio tra ripassare dopo tre
giorni o dopo cinque, per un esame lontano, scegli cinque. L'errore che
fanno quasi tutti gli studenti non è aspettare troppo: è ripassare
troppo presto, quando il ricordo è ancora così fresco che la ripresa non
serve quasi a nulla, e poi non tornarci più.

Resta un dubbio pratico. Le applicazioni di flashcard e molti manuali
di metodo di studio consigliano intervalli \emph{crescenti}: ripassare
dopo un giorno, poi dopo tre, poi dopo una settimana, poi dopo un mese.
È davvero meglio che ripassare a intervalli regolari? Alice Latimier,
Hugo Peyre e Franck Ramus hanno raccolto tutti gli studi che avevano
confrontato le due soluzioni, nel caso in cui ogni ripresa sia una
pratica del recupero \parencite{latimier2021}. Il risultato è doppio.
Da una parte, distribuire il recupero invece di concentrarlo produce un
vantaggio forte, di quelli che nella ricerca educativa si incontrano di
rado. Dall'altra, tra intervalli crescenti e intervalli uniformi non
c'è una differenza apprezzabile. Conta distanziare; il calendario
preciso conta molto meno di quanto si creda. È una buona notizia,
perché ti libera dall'ansia del programma perfetto: un calendario
ragionevole, rispettato, vale più di un calendario ottimo, abbandonato
dopo una settimana.

Anche qui, come per il recupero, la matematica fa un po' storia a sé.
Una meta-analisi recente trova che distribuire gli esercizi di
matematica nel tempo aiuta, e aiuta anche in corsi reali, ma con un
effetto più piccolo che in altri campi \parencite{murray2025}. Più
piccolo non vuol dire trascurabile: vuol dire che per un esame di
analisi o di statistica la distribuzione è una condizione necessaria,
non una bacchetta magica, e che va accompagnata dall'alternanza dei
tipi di problemi di cui parleremo nel capitolo~\ref{cap:alternanza}.

\begin{table}[tbp]
  \centering\small
  \begin{tabular}{@{}>{\raggedright\arraybackslash}p{0.34\textwidth}>{\raggedright\arraybackslash}p{0.56\textwidth}@{}}
    \toprule
    \textit{Devi ricordare per\dots} & \textit{Intervallo indicativo tra le riprese} \\
    \midrule
    una settimana (una prova in itinere) & 1--2 giorni \\
    \addlinespace
    un mese & circa una settimana \\
    \addlinespace
    tre mesi (un esame a fine semestre) & 2--3 settimane \\
    \addlinespace
    un anno (un concorso, un esame di Stato) & 2--4 settimane \\
    \addlinespace
    per sempre & riprese sempre più rade, ogni pochi mesi \\
    \bottomrule
  \end{tabular}
  \caption{Quanto lasciar passare tra una ripresa e l'altra. Valori
    orientativi ricavati dalla \enquote{cresta} di
    \textcite{cepeda2008}: l'intervallo migliore è circa un quinto del
    tempo per cui si vuole ricordare, se si tratta di poche settimane,
    e circa un ventesimo, se si tratta di un anno. Nel dubbio, meglio
    un intervallo un po' più lungo che uno più breve.}
  \label{tab:intervalli}
\end{table}

\section{Il ripasso a successive riprese}

Fin qui abbiamo parlato di \emph{quando} tornare su ciò che hai
studiato. Ma che cosa fare, quando ci torni? Se la risposta è
\enquote{rileggere}, buona parte del vantaggio va perduta: la
rilettura distribuita è meglio della rilettura concentrata, ma resta
una rilettura (capitolo~\ref{cap:inganni}). Il meglio si ottiene
quando le due forze di questo capitolo e del precedente lavorano
insieme: \emph{recuperare}, e farlo \emph{a intervalli}.

Katherine Rawson e John Dunlosky, della Kent State University, hanno
dato a questa combinazione una forma precisa, e un nome: \emph{ripasso
a successive riprese} (\emph{successive relearning})
\parencite{rawson2013, rawson2022}. La procedura ha due regole. La
prima: in ogni sessione, ti interroghi sui concetti da imparare -- con
domande o flashcard, a libro chiuso -- e ogni volta che sbagli controlli
la risposta e rimetti la domanda in fondo al mazzo, finché non hai
risposto correttamente a tutte. Questo è il \enquote{criterio}: non si
chiude la sessione finché ogni domanda non ha avuto almeno una risposta
giusta. La seconda regola: la stessa procedura si ripete in più
sessioni distanziate nel tempo, di solito da tre a cinque, a distanza
di giorni o di settimane.

Nei loro esperimenti, condotti con studenti universitari in corsi veri
e su argomenti che sarebbero stati chiesti all'esame, il ripasso a
successive riprese ha prodotto miglioramenti consistenti nei voti degli
esami del corso e, soprattutto, un ricordo che a distanza di mesi era
molto migliore di quello degli studenti che avevano studiato nel modo
abituale. C'è poi un dettaglio
che rende la procedura molto più sostenibile di quanto sembri: le
sessioni successive alla prima sono \emph{brevi}. La prima volta ci
vuole tempo per arrivare al criterio, perché molte risposte sono
sbagliate; la seconda ne servono meno; alla terza e alla quarta, per lo
stesso materiale, bastano pochi minuti. Il tempo totale non è molto
superiore a quello di chi studia nel modo abituale. Cambia la sua
distribuzione.

Nel capitolo precedente abbiamo visto un errore comune: eliminare una
domanda dal mazzo dopo la prima risposta giusta. Il ripasso a
successive riprese spiega perché è un errore e, insieme, come
rimediare. Dentro una singola sessione, una risposta giusta basta:
puoi mettere la scheda da parte. Ma non per sempre, soltanto fino alla
sessione successiva, tra qualche giorno. Una domanda esce davvero dal
mazzo solo quando hai risposto correttamente in più sessioni, in giorni
diversi, con intervalli che si sono via via allungati. A quel punto non
la sai \enquote{ancora}: la sai.

Se dovessi scegliere una sola procedura di studio, una sola, da portare
con te per tutta l'università, dovrebbe essere questa. Contiene quasi
tutto ciò che la ricerca sa sull'apprendimento duraturo: il recupero,
la distribuzione, il controllo delle risposte, l'attenzione concentrata
su ciò che non si sa ancora. E non richiede nulla di speciale: un mazzo
di schede, o un foglio di domande, e un calendario.

\section{Costruire un calendario di ripasso}

Il principio è chiaro; la difficoltà è pratica. Distribuire lo studio
vuol dire decidere \emph{oggi} che cosa farai tra tre giorni e tra tre
settimane, e poi farlo davvero, quando nessuno te lo chiede e magari
hai altre scadenze più vicine. Senza un calendario scritto, la
distribuzione non avviene: le buone intenzioni cedono alla scadenza
più urgente, e si torna, quasi senza accorgersene, allo studio della
vigilia.

C'è anche un ostacolo che non dipende da te. La scuola italiana, e
ancora di più l'università, sono organizzate per unità chiuse:
\enquote{finito il Risorgimento, si passa all'Unità}; finito il corso,
si dà l'esame; dato l'esame, si passa al successivo. La struttura
tipica di un insegnamento universitario -- lezioni durante il semestre,
studio concentrato nelle settimane della sessione, esame, oblio -- è
quasi la definizione di pratica concentrata. Nessuno ti obbliga a
riprendere ciò che hai studiato a ottobre prima di gennaio, né ciò che
hai studiato per l'esame di gennaio dopo averlo superato. Se vuoi
distribuire, devi farlo tu. Vediamo tre situazioni tipiche.

\subsection*{Una prova in itinere tra una settimana}

Mettiamo che tra sette giorni tu abbia una prova intermedia sulle
prime quattro settimane di un corso. Il tempo è poco, ma sufficiente
per tre riprese. Il primo giorno studi il materiale per capirlo e, alla
fine, fai il primo recupero a libro chiuso (il ciclo in tre passi del
capitolo~\ref{cap:recupero}); prepari le domande sui punti essenziali.
Il secondo giorno rispondi alle domande senza guardare, fino al
criterio, e ristudi solo ciò che hai sbagliato. Il quarto o quinto
giorno ripeti; il sesto fai una simulazione della prova, nelle
condizioni della prova. Il giorno prima: una breve ripresa, e a letto
presto. Sono quattro o cinque sessioni, la maggior parte brevi, al
posto di una lunga giornata alla vigilia.

\subsection*{Un esame a fine semestre}

Qui il tempo è il tuo migliore alleato, a condizione di cominciare
subito. La regola è semplice: \emph{pianifica il semestre, non la
sessione}. Per ogni insegnamento, fissa un appuntamento settimanale di
recupero -- mezz'ora o un'ora, sempre lo stesso giorno -- in cui ti
interroghi non solo sugli argomenti della settimana, ma anche su quelli
delle settimane precedenti, via via più rade. Un argomento spiegato alla
prima settimana di lezione verrà così ripreso, nel corso del semestre,
cinque o sei volte, a intervalli che si allungano; quando arriverà la
sessione, non dovrai \emph{studiarlo}, ma soltanto \emph{ripassarlo}.
Nelle ultime due settimane prima dell'appello il lavoro cambia
natura: niente materiale nuovo, simulazioni d'esame, ripresa di tutto
il programma per blocchi mescolati. Ne parleremo in dettaglio nel
capitolo~\ref{cap:metodo} e nel capitolo~\ref{cap:semestre}.

\subsection*{Un concorso o un esame di Stato tra un anno}

Quando l'orizzonte è lungo e il programma sterminato -- un concorso
pubblico, l'esame di abilitazione, il test di specializzazione in
medicina -- la distribuzione non è un'opzione: è l'unico modo di
tenere tutto in memoria nello stesso momento. Dividi il programma in
blocchi e studiali uno dopo l'altro; ma dedica ogni settimana una parte
fissa del tempo, per esempio un quarto, a riprendere con il recupero i
blocchi già studiati. Per un orizzonte di un anno, gli intervalli tra
le riprese di uno stesso blocco possono arrivare a due, tre o quattro
settimane (tabella~\ref{tab:intervalli}). All'inizio la parte di
ripresa ti sembrerà tempo sottratto al programma nuovo; verso la fine
ti accorgerai che è l'unica ragione per cui ricordi ancora il primo
blocco.

\subsection*{Carta, quaderno o applicazione}

Per organizzare le riprese servono due strumenti: un calendario e un
archivio di domande. Il calendario può essere un'agenda di carta o
quello del telefono; l'importante è che le riprese siano scritte come
appuntamenti, con un giorno e un'ora, e non come vaghi propositi.
L'archivio di domande può essere un quaderno con le domande a sinistra e
le risposte a destra, un mazzo di schede di carta divise in scatole
(chi risponde bene sposta la scheda nella scatola successiva, che si
riprende più di rado; chi sbaglia la riporta nella prima), oppure una
delle molte applicazioni di flashcard che calcolano da sole quando
riproporre ciascuna domanda. Le applicazioni sono comode e fanno bene
il loro mestiere, purché tu le usi come si deve: rispondendo prima di
girare la scheda, e valutando onestamente la tua risposta. Ma ricorda
il risultato di Latimier: l'algoritmo perfetto conta meno della
costanza. Una scatola da scarpe usata ogni giorno vale più di
un'applicazione sofisticata aperta una volta al mese.

\begin{inpratica}
\textbf{Il calendario di ripasso in quattro mosse.}
\begin{enumerate}
  \item \emph{Segna le date}: esami, prove in itinere, consegne. Per
        ciascuna, guarda quanto tempo manca e scegli l'intervallo tra
        le riprese (tabella~\ref{tab:intervalli}).
  \item \emph{Fissa gli appuntamenti}: per ogni insegnamento, una
        sessione di recupero alla settimana, sempre lo stesso giorno,
        scritta in agenda come una lezione.
  \item \emph{Riprendi il vecchio insieme al nuovo}: in ogni sessione,
        domande sugli argomenti della settimana e su quelli delle
        settimane precedenti, fino al criterio.
  \item \emph{Allunga gli intervalli}: ciò che sai bene torna più di
        rado; ciò che sbagli torna presto.
\end{enumerate}
Un modello di calendario da fotocopiare è nell'appendice~\ref{app:schede}.
\end{inpratica}

\section{Contro la notte prima dell'esame}

Se la distribuzione funziona così bene, perché quasi tutti gli
studenti continuano a concentrare lo studio nei giorni che precedono
l'esame? Le ragioni esterne le abbiamo viste: il calendario delle
sessioni, gli esami che si accumulano, le scadenze più vicine che
schiacciano quelle più lontane. Ma c'è anche una ragione interna, più
insidiosa: lo studio concentrato \emph{sembra} funzionare meglio.

Nate Kornell lo ha mostrato con un esperimento semplice
\parencite{kornell2009}. Alcuni studenti imparavano con le flashcard
vocaboli inglesi difficili, con la loro definizione; una parte delle
schede veniva ripetuta a intervalli brevissimi, in piccoli mazzi
ripresi più volte di seguito, un'altra a intervalli più lunghi. Al test finale la
distribuzione si è rivelata più efficace per 63 partecipanti su 70:
circa nove su dieci. Eppure, quando è stato chiesto loro quale metodo
avesse funzionato meglio, 42 dei 58 che hanno espresso un giudizio --
quasi tre su quattro -- hanno indicato lo studio concentrato. Hanno
imparato di più in un modo e hanno creduto di aver imparato di più
nell'altro.

La spiegazione è quella che conosci già dal
capitolo~\ref{cap:inganni}. Mentre studi in modo concentrato, tutto è
fresco: ogni informazione è ancora lì, a portata di mano; le risposte
arrivano pronte; la sensazione di padronanza è forte. Quello che
misuri, in quel momento, è la forza di recupero, non l'apprendimento.
Nello studio distribuito accade il contrario: ogni volta che riprendi
il materiale, una parte l'hai dimenticata, e la sensazione è quella di
dover ricominciare da capo. È una sensazione sgradevole, e ingannevole.
L'oblio parziale che la provoca è esattamente ciò che rende la ripresa
efficace.

Lo studio della vigilia, poi, \enquote{funziona} davvero, in un senso
preciso: se l'esame è domani e ti chiede di riconoscere o di ripetere
ciò che hai letto oggi, puoi superarlo. Il prezzo si paga dopo. Ciò
che è stato stipato in pochi giorni se ne va in pochi giorni; quando il
corso dell'anno successivo lo presuppone, bisogna ristudiarlo quasi da
capo. È il secondo dei tre errori del capitolo~\ref{cap:vita}:
studiare per il giorno dopo, non per la vita.

E c'è un aggravante, quando la vigilia diventa una notte in bianco.
Il sonno non è una pausa per la memoria: è una delle fasi in cui la
memoria lavora di più. Già nel 1924 John Jenkins e Karl Dallenbach
avevano osservato che si dimentica meno dopo alcune ore di sonno che
dopo lo stesso tempo passato svegli \parencite{jenkins1924}; oggi
sappiamo che durante il sonno i ricordi della giornata vengono
riattivati e consolidati. Chi rinuncia a dormire per studiare di più
somma due svantaggi: studia in modo concentrato e toglie a ciò che ha
studiato la notte che dovrebbe fissarlo. E se la notte in bianco cade
prima di una giornata di studio, compromette anche la capacità di
imparare cose nuove il giorno dopo \parencite{yoo2007}. Ne parleremo nel
capitolo~\ref{cap:corpo}; per ora basta una regola: \emph{la notte
prima dell'esame si dorme}.

\begin{riquadro}{Per chi insegna}
Il corso universitario tipico è una sequenza di unità chiuse, e
l'esame di ciascuna sembra dire allo studente che, superata la prova,
quell'unità si può dimenticare. Chi insegna può cambiare questo
messaggio senza aggiungere una sola lezione: basta che i quiz, le
prove in itinere e l'esame finale contengano anche domande sugli
argomenti delle settimane precedenti, e non soltanto su quelli
dell'ultima. Una prova \emph{cumulativa} obbliga a riprendere, a
intervalli che si allungano, ciò che è stato già spiegato: è il
ripasso a successive riprese, imposto dal calendario invece che
lasciato alla buona volontà di ciascuno. Va detto agli studenti
all'inizio del corso, perché la scelgano invece di subirla
(capitoli~\ref{cap:docenti} e~\ref{cap:valutare}).
\end{riquadro}

\asterismo

Io credo che il messaggio di questo capitolo sia, in fondo, un
messaggio di sollievo. Non serve studiare di più: spesso serve
studiare lo stesso tempo, ma distribuito diversamente. Un'ora al giorno
per tre giorni vale più di tre ore in un giorno; mezz'ora alla
settimana per tutto il semestre vale più di una settimana di clausura
prima dell'appello. Quintiliano chiedeva un esercizio \enquote{se
possibile ogni giorno}; Ovidio vedeva la goccia che scava la pietra. Il
tempo non è soltanto ciò che ti fa dimenticare. Se lo usi bene, è ciò
che ti fa ricordare.

\begin{daricordare}
Il tempo è un alleato, non un nemico. Distribuisci lo studio in più
sessioni invece di concentrarlo, e in ogni ripresa interrogati invece
di rileggere. Più a lungo vuoi ricordare, più distanti possono essere
le riprese.
\end{daricordare}

\begin{domande}
  \item Perché tre ore di studio il giorno prima dell'esame rendono
        meno di un'ora al giorno per tre giorni? Rispondi usando le
        due forze della memoria.
  \item Che cos'è la \enquote{cresta} scoperta da Cepeda e colleghi?
        Quanto conviene aspettare prima di riprendere un argomento per
        un esame tra tre mesi?
  \item Descrivi il ripasso a successive riprese: quali sono le sue due
        regole?
  \item Scrivi il calendario delle riprese per il prossimo esame che
        devi sostenere.
  \item (Ripresa, capitolo~\ref{cap:recupero}) Qual è l'errore più
        comune con le flashcard, e come lo corregge il ripasso a
        successive riprese?
  \item (Ripresa, capitolo~\ref{cap:vita}) Che cosa aveva misurato
        Ebbinghaus confrontando 68 e 38 ripetizioni?
\end{domande}

\begin{sintesi}
  \item A parità di tempo, distribuire lo studio in più sessioni
        produce un ricordo molto più duraturo che concentrarlo: è uno
        dei risultati più replicati della psicologia.
  \item Funziona perché un po' di oblio tra una ripresa e l'altra
        rende più faticoso, e quindi più efficace, il ricordo
        successivo.
  \item L'intervallo migliore tra le riprese è una frazione del tempo
        per cui si vuole ricordare; intervalli troppo brevi costano più
        di intervalli un po' troppo lunghi, e il calendario preciso
        conta meno della costanza.
  \item Il ripasso a successive riprese unisce recupero e
        distribuzione: interrogarsi fino al criterio, in tre-cinque
        sessioni distanziate, sempre più brevi.
  \item Lo studio della vigilia sembra funzionare perché misura la
        forza di recupero, non l'apprendimento; la notte in bianco
        aggiunge al danno la perdita del sonno.
\end{sintesi}

\finecapitolo
